% Supplementary Material
\section{Supplementary Material}
\label{sec:supplementary}

This appendix provides extended discussion, detailed examples, and supporting analysis
referenced from the main text. Sections are ordered to follow the main paper's
structure.

% Ordering follows main text reference order:
% A.1 Threshold Design Rationale (from Section 4)
% A.2 Complete Metric Rubrics (from Section 4)
% A.3 Implementation Details (from Section 5)
% A.4 Integration Considerations (from Section 5)
% A.5 Detailed Case Studies (from Section 6)
% A.6 Failure Modes and Mitigations (from Section 6)
% A.7 Evaluation Tables (from Section 6)
% A.8 Cost Analysis and ROI Framework (from Section 6)
% A.9 When to Use SPOQ (from Section 7)

\subsection{Threshold Design Rationale}
\label{subsec:threshold-rationale}

The 95/90 (planning) and 95/80 (code) thresholds were chosen based on
practical experience:

\begin{itemize}
    \item \textbf{95\% aggregate} ensures overall high quality while allowing
    minor imperfections in individual metrics.

    \item \textbf{90\% per-metric for planning} reflects that planning errors
    propagate downstream. A plan with 70\% dependency graph quality will
    cause execution failures regardless of other metrics.

    \item \textbf{80\% per-metric for code} acknowledges legitimate tradeoffs.
    A task might score 75\% on SOLID adherence for pragmatic reasons while
    still being acceptable.

    \item \textbf{Plans are cheap to fix.} Re-decomposing tasks costs human
    time but no agent compute. Code rework costs both.
\end{itemize}

\subsection{Complete Metric Rubrics}
\label{sec:appendix-metrics}

This subsection provides detailed scoring rubrics for both validation gates.

\subsubsection{Planning Validation Rubrics}

\paragraph{Vision Clarity (VC).}
\begin{itemize}
    \item 100: Clear overview with problem statement, solution, scope boundaries,
    and end state
    \item 80: Clear overview, minor ambiguity in scope
    \item 60: Overview present but vague
    \item 40: Overview missing key context
    \item 0: No overview or incomprehensible goal
\end{itemize}

\paragraph{Architecture Quality (AQ).}
\begin{itemize}
    \item 100: ASCII/visual diagram with all components, relationships, and data flow
    \item 80: Diagram present, minor gaps in explanation
    \item 60: Text-only architecture, no diagram
    \item 40: Partial architecture description
    \item 0: No architecture section
\end{itemize}

\paragraph{Task Decomposition (TD).}
\begin{itemize}
    \item 100: Atomic tasks (1--4h), independent where possible, complete coverage
    \item 80: Good decomposition, minor overlap
    \item 60: Some tasks too large or overlapping
    \item 40: Significant decomposition issues
    \item 0: Tasks not properly decomposed
\end{itemize}

\paragraph{Dependency Graph (DG).}
\begin{itemize}
    \item 100: Visual graph, all dependencies valid, no cycles, optimal ordering
    \item 80: Graph present, dependencies valid, minor ordering improvements possible
    \item 60: Text-only dependencies, all valid
    \item 40: Some invalid dependencies or missing graph
    \item 0: Circular dependencies or critically broken graph
\end{itemize}

\paragraph{Coverage Completeness (CC).}
Score equals percentage of success criteria mapped to tasks:
$\text{CC} = 100 \times \frac{|\text{covered criteria}|}{|\text{total criteria}|}$

\paragraph{Phase Ordering (PO).}
\begin{itemize}
    \item 100: Phases follow logical progression, maximum parallelism exploited
    \item 90: Good ordering, minor parallelism opportunities missed
    \item 75: Ordering works but inefficient
    \item 60: Some ordering violations
    \item 0: Critical ordering errors (dependent tasks in same wave)
\end{itemize}

\paragraph{Scope Coherence (SC).}
\begin{itemize}
    \item 100: All tasks directly serve epic goal
    \item 80: 1 tangential task
    \item 60: 2--3 tangential tasks
    \item 40: Multiple unrelated tasks
    \item 0: Tasks don't align with epic goal
\end{itemize}

\paragraph{Success Criteria Quality (SQ).}
\begin{itemize}
    \item 100: All criteria SMART, checkbox format, testable
    \item 90: Criteria measurable, minor gaps
    \item 75: Some criteria vague
    \item 60: Multiple unmeasurable criteria
    \item 0: No success criteria or all vague
\end{itemize}

\paragraph{Risk Identification (RI).}
\begin{itemize}
    \item 100: Risks section with likelihood, impact, and mitigations
    \item 90: Risks mentioned, implicit mitigations
    \item 75: Some risks noted, no mitigations
    \item 60: Risks not addressed
    \item 0: Critical risks ignored
\end{itemize}

\paragraph{Integration Strategy (IS).}
\begin{itemize}
    \item 100: Clear integration points, verification steps between phases
    \item 90: Integration implicit but clear
    \item 75: Some integration points unclear
    \item 60: Integration strategy missing
    \item 0: Tasks cannot be integrated as designed
\end{itemize}

\subsubsection{Code Validation Rubrics}

\paragraph{Syntactic Correctness (SC).}
\begin{itemize}
    \item 100: Compiles cleanly, 0 warnings
    \item 80: Compiles with minor warnings
    \item 60: Compiles with significant warnings
    \item 0: Does not compile
\end{itemize}

\paragraph{Test Existence (TE).}
Score equals percentage of new public methods with corresponding tests:
$\text{TE} = 100 \times \frac{|\text{tested methods}|}{|\text{new public methods}|}$

\paragraph{Test Pass Rate (TP).}
Score equals percentage of tests passing:
$\text{TP} = 100 \times \frac{|\text{passing tests}|}{|\text{total tests}|}$

\paragraph{Requirements Fidelity (RF).}
\begin{itemize}
    \item 100: All requirements fully implemented
    \item 80: All core requirements met, minor gaps
    \item 60: Core requirements met, some missing
    \item 40: Partial implementation
    \item 0: Does not address requirements
\end{itemize}

\paragraph{SOLID Adherence (SA).}
20 points per principle (S, O, L, I, D) based on degree of adherence.

\paragraph{Security (SE).}
Start at 100, deduct for OWASP Top 10 vulnerabilities:
SQL injection ($-100$), command injection ($-100$), XSS ($-60$),
broken auth ($-80$), sensitive data exposure ($-60$), etc.

\paragraph{Error Handling (EH).}
20 points each for: I/O wrapped in try-catch, meaningful error messages,
proper logging, resource cleanup, safe user-facing errors.

\paragraph{Scalability (SL).}
Based on algorithm complexity of hot paths:
O(1)/O(log n)/O(n) = 100, O(n log n) = 90, O(n$^2$) = 40, O(2$^n$) = 0.

\paragraph{Code Clarity (CC).}
\begin{itemize}
    \item 100: Crystal clear, reads like well-written prose
    \item 80: Clear with minor naming improvements possible
    \item 60: Understandable but requires effort
    \item 40: Confusing structure or naming
    \item 0: Unreadable, magic numbers, cryptic names
\end{itemize}

\paragraph{Completeness (CO).}
\begin{itemize}
    \item 100: Complete, no loose ends
    \item 80: Minor polish needed
    \item 60: Some TODOs remain but core is done
    \item 40: Significant unfinished sections
    \item 0: Stub implementations, placeholders
\end{itemize}

Automatic deductions: TODO ($-25$), FIXME ($-25$), NotImplementedException ($-30$).

\subsection{Implementation Details}
\label{subsec:implementation-detail}

This subsection provides complete schema examples, journal entry formatting, and
skill metadata referenced from Section~\ref{sec:implementation}.

\subsubsection{Epic Directory Structure}

Each epic occupies its own directory under \texttt{spoq/epics/}:

\begin{lstlisting}[language=bash,caption={Epic directory layout}]
spoq/epics/
  active/                   # Epics in progress
    epic-name/
      EPIC.md
      tasks/
        01-init-project.yml
        ...
  complete/                 # Finished epics
  ROADMAP.md                # Priority tracker
\end{lstlisting}

The \texttt{EPIC.md} file provides context: goal statement, architecture diagrams,
success criteria, dependency visualization, wave assignments, effort estimates,
and risk assessment. Individual task files contain execution-ready specifications.

\subsubsection{Task YAML Schema Examples}

The task specification (Definition~\ref{def:task-spec}) organizes fields into
three categories. The following listings illustrate each category with a
representative task.

\paragraph{Identity Fields.}
These establish task context within an epic:

\begin{lstlisting}[,caption={Task identity fields}]
id: 04-content-constants       # Unique within epic
title: Create Content Constants File
epic: spoq-website             # Parent epic reference
\end{lstlisting}

\needspace{10cm}
\paragraph{Execution Control Fields.}
These govern scheduling, effort estimation, and dependency resolution:

\begin{lstlisting}[,caption={Execution control fields}]
status: pending                # pending | in_progress | completed
priority: high                 # critical | high | medium | low
phase: 1                       # Wave assignment (0 = no deps)

estimate:                      # PERT three-point estimate
  optimistic: 15m
  realistic: 45m
  pessimistic: 2h

dependencies:                  # Task IDs that must complete first
  - 01-init-project
  - 02-setup-deps

skills_required:               # Domain knowledge needed
  - typescript
  - react
\end{lstlisting}

\paragraph{Deliverable and Verification Fields.}
These define expected outputs and acceptance criteria:

\begin{lstlisting}[,caption={Deliverable and verification fields}]
files_to_touch:                # All files to be modified
  - src/lib/constants.ts
  - tests/constants.test.ts

outputs:                       # Tangible deliverables
  - "Constants file with typed exports"
  - "Unit tests achieving 100% coverage"

acceptance_criteria:           # Verification checklist
  - "[ ] TypeScript compiles without errors"
  - "[ ] All tests pass: `npm test constants`"
  - "[ ] No hardcoded strings in component files"
\end{lstlisting}

\needspace{12cm}
\paragraph{Description Field.}
The description provides structured implementation guidance using Markdown:

\begin{lstlisting}[,caption={Task description structure}]
description: |
  ## Objective
  Create a centralized constants file for UI strings.

  ## Steps
  1. Create src/lib/constants.ts with typed exports
     ```typescript
     export const SITE_NAME = "SPOQ" as const;
     ```
  2. Add unit tests in tests/constants.test.ts
  3. Replace hardcoded strings in existing components

  ## Verification
  ```bash
  npm run typecheck && npm test constants
  ```
\end{lstlisting}

\needspace{4cm}
\subsubsection{Journal Entry Format}

Journal entries (Definition~\ref{def:journal-entry}) use YAML frontmatter
followed by Markdown sections. The frontmatter captures machine-readable
metadata:

\begin{lstlisting}[,caption={Journal entry frontmatter}]
---
agent: Claude Code (Opus 4.5)
start_time: 2025-11-01T10:52:34Z  # ISO 8601 UTC
end_time: 2025-11-01T11:23:15Z
confidence: 0.92                  # 0.0-1.0 calibration score
session_type: development         # development | refactor | bugfix | ...
files_modified:
  - src/components/Hero.tsx
  - tests/hero.test.tsx
tasks_completed: 2
tasks_total: 3
---
\end{lstlisting}

\needspace{14cm}
The body follows a standardized section layout for consistent parsing:

\begin{lstlisting}[,caption={Journal entry body structure}]
## Summary
Brief 1-2 sentence overview of accomplishments.

## Work Completed
- Task 09: Hero section component
- Task 10: Features grid with responsive layout

## Changes Made
**Frontend Components**
- `Hero.tsx` - Implemented animated gradient background
- `Features.tsx` - Added 6-card responsive grid

## Issues Encountered
None (or specific issues with resolutions)

## Testing
- Unit tests: PASS (12/12)
- Integration tests: PASS (4/4)

## Next Steps
1. Implement call-to-action section
2. Add accessibility attributes
\end{lstlisting}

\paragraph{Tooling Support.}
SPOQ provides utility scripts for journal management:
\begin{itemize}
    \item \texttt{get-time.py}: Captures accurate UTC timestamps
    \item \texttt{archive-journal.py}: Auto-archives when exceeding 1500 lines
    \item \texttt{parse-to-db.py}: Extracts entries into SQLite for analysis
    \item \texttt{create-mega-journal.py}: Consolidates archives for reporting
\end{itemize}

\subsubsection{Confidence Score Interpretation}

\begin{table}[htbp]
\centering
\caption{Confidence Score Interpretation}
\label{tab:confidence}
\small
\begin{tabular}{ll}
\toprule
\textbf{Score Range} & \textbf{Interpretation} \\
\midrule
0.95--1.0 & Thoroughly tested, production-ready \\
0.85--0.94 & Well tested, minor edge cases possible \\
0.75--0.84 & Functional, additional testing recommended \\
0.65--0.74 & Works but requires validation \\
$<$0.65 & Experimental or known issues present \\
\bottomrule
\end{tabular}
\end{table}

\needspace{4cm}
\subsubsection{Skill Inventory}

\begin{table}[htbp]
\centering
\caption{SPOQ Skill Inventory}
\label{tab:skills}
\small
\begin{tabular}{llp{5cm}}
\toprule
\textbf{Skill} & \textbf{Stage} & \textbf{Purpose} \\
\midrule
epic-planning & Planning & Decompose goals into epics with dependency graphs \\
epic-validation & Validation & Score epics against 10 planning metrics \\
task-validation & Validation & Score individual tasks before execution \\
agent-execution & Execution & Orchestrate parallel agent swarms \\
agent-validation & Validation & Score completed work against 10 code metrics \\
journal-tracker & Cross-cutting & Track sessions with confidence scores \\
\bottomrule
\end{tabular}
\end{table}

\subsubsection{Skill Metadata and Invocation}

Skills are invoked via slash commands that expand into full prompts:
\begin{lstlisting}[language=bash,caption={Skill invocation examples}]
# Plan a new epic
/epic-planning "Implement user authentication system"

# Validate before execution
/epic-validation @spoq/epics/active/auth-system

# Execute with parallel agents
/agent-execution @spoq/epics/active/auth-system

# Validate completed work
/agent-validation
\end{lstlisting}

Each skill's \texttt{SKILL.md} includes YAML frontmatter specifying metadata:

\needspace{6cm}
\begin{lstlisting}[,caption={Skill metadata format}]
---
name: epic-planning
description: Decompose high-level goals into structured epics
             with atomic tasks and dependency DAGs.
---
\end{lstlisting}

The body provides:
\begin{itemize}
    \item \textbf{When to Use:} Activation criteria
    \item \textbf{Core Patterns:} Documented approaches with examples
    \item \textbf{Quality Criteria:} Verification standards
    \item \textbf{Integration Points:} Connections to other skills
\end{itemize}

\needspace{6cm}
\subsection{Integration Considerations}
\label{subsec:integration-detail}
\label{subsec:tooling-integration}

A recurring question from practitioners concerns SPOQ's integration
with existing development infrastructure. We address this honestly:
SPOQ currently operates as a standalone orchestration layer invoked
via slash commands, without native integrations with external tools.
However, its design accommodates several integration patterns.

\paragraph{CI/CD Pipeline Patterns.}
SPOQ's wave-based execution maps naturally to CI/CD pipeline stages.
Wave 0 corresponds to setup and dependency installation, middle waves
to parallel build and test jobs, and final waves to integration testing
and deployment. In our infrastructure case study, the GitLab CI pipeline
was updated as a task within the SPOQ epic itself, demonstrating that
CI/CD configuration is orchestrable work rather than requiring special
integration.

For automated triggering, organizations could invoke SPOQ from pipeline
scripts when epic YAML files change, treating the orchestrator as a
build step. A GitLab CI job might monitor the \texttt{spoq/epics/}
directory and trigger agent execution on detected changes. This pattern
remains unexplored in our evaluation but represents a natural extension.
GitHub Actions workflows could similarly invoke SPOQ as a step,
potentially enabling fully automated epic execution on pull request events.

\paragraph{Project Management Alignment.}
SPOQ task YAML files parallel the structure of Jira stories and Linear
issues: both contain titles, descriptions, acceptance criteria, and
estimates. The structural similarity suggests bidirectional sync is
feasible; a task defined in SPOQ could create a corresponding Jira
ticket, and vice versa. However, SPOQ does not currently implement
such synchronization, which may create duplicate tracking overhead
for teams already committed to existing PM tools.

The journal system partially addresses visibility by providing an
audit trail of work completed, though it is designed for agent
coordination rather than project management. Teams could export
journal entries to their PM tools as a manual bridge, accepting
the overhead until native integration warrants development investment.

\paragraph{Git Workflow Recommendations.}
Parallel agent execution raises questions about branch strategy.
Our current approach uses a single working branch with agents
operating on disjoint files, ensured by task \texttt{files\_to\_touch}
specifications. This works well for up to 12 concurrent agents
but may encounter contention at larger scales.

For larger teams or epics with overlapping file modifications,
we recommend git subtrees for epic isolation: each epic operates
in its own subtree, with automated merges after validation gates
pass. Feature branches per wave provide an alternative, with
wave N's branch merging into main before wave N+1 begins. Both
approaches add orchestration complexity but reduce merge conflicts.

\paragraph{IDE Integration Possibilities.}
SPOQ's slash-command interface suggests IDE integration opportunities.
A VS Code extension could display the current epic's dependency graph,
highlight files assigned to each task, and show validation scores
in real time. JetBrains plugins could integrate journal entries
into the project view, surfacing agent work history alongside
version control logs.

Such integrations remain future work. Our current implementation
prioritizes methodology validation over tooling polish, reflecting
a conscious choice to prove the approach before investing in
developer experience enhancements.

\paragraph{Future Integration Directions.}
Several integration directions merit exploration:
(1) native Jira/Linear/Asana adapters for bidirectional task sync;
(2) Slack/Teams webhooks for wave completion notifications;
(3) GitHub/GitLab API integration for automated PR creation per wave;
(4) IDE extensions for real-time epic visualization.
These would reduce friction for teams adopting SPOQ within established
toolchains, though each adds maintenance burden. We advocate starting
with SPOQ as a standalone layer, adding integrations only when
adoption justifies the investment.

\subsection{Detailed Case Studies}
\label{subsec:case-study-detail}

\subsubsection{Case Study 1: UI Improvements Epic}
\label{subsubsec:ui-epic-detail}

\paragraph{Epic Overview.}
The UI improvements epic modernized a monitoring dashboard with toast
notifications, data tables, and API key management components. The
epic comprised 13 tasks across 2 waves.

\paragraph{Wave Structure.}
\begin{itemize}
    \item \textbf{Wave 0 (12 tasks):} Independent component implementations
    including Sonner setup, DataTable component, modal dialogs, toast
    integrations, search functionality, and tests.
    \item \textbf{Wave 1 (1 task):} End-to-end QA depending on all Wave 0 tasks.
\end{itemize}

\paragraph{Execution Metrics.}
Table~\ref{tab:ui-metrics} summarizes the execution.

\begin{table}[htbp]
\centering
\caption{UI Improvements Epic Execution}
\label{tab:ui-metrics}
\begin{tabular}{ll}
\toprule
\textbf{Metric} & \textbf{Value} \\
\midrule
Total tasks & 13 \\
Wave 0 parallelism & 12 concurrent agents \\
Wave 1 parallelism & 1 agent \\
Sequential estimate & 18.5 hours \\
Parallel estimate & 3.5 hours \\
Speedup factor & 5.3$\times$ \\
Tasks completed & 12 of 13 (92\%) \\
Rework cycles & 2 (tasks 04, 07) \\
\bottomrule
\end{tabular}
\end{table}

\paragraph{Failure Mode Encountered.}
Task 04 entered a runaway retry loop, executing \texttt{npm install sonner}
over 100 times. The agent failed to recognize that the dependency was
already installed and continued retrying indefinitely.

\paragraph{Lesson Learned.}
SPOQ now enforces a maximum of 3 retries per installation command and
requires agents to verify existing dependencies before installation attempts.

\subsubsection{Case Study 2: Client Website Rebrand}
\label{subsubsec:rebrand-epic-detail}

\paragraph{Epic Overview.}
This epic rebranded an external client's sales website (Pinpoint, a B2B QA
testing platform) from founder-centric messaging to developer-focused
positioning. The work included removing biographical content, rewriting
homepage copy, creating three persona-specific landing pages, updating
navigation and SEO metadata, and expanding the test suite. The epic comprised
12 tasks across 4 waves and was executed on a codebase maintained by a
separate engineering team.

\paragraph{Wave Structure.}
\begin{itemize}
    \item \textbf{Wave 0 (2 tasks):} Content removal and routing scaffold
    (parallel).
    \item \textbf{Wave 1 (5 tasks):} Homepage rewrite, pricing update, and three
    persona landing pages (parallel).
    \item \textbf{Wave 2 (3 tasks):} Navigation, SEO, and section refinement
    (parallel).
    \item \textbf{Wave 3 (2 tasks):} Test expansion and accessibility polish
    (parallel).
\end{itemize}

\paragraph{Execution Metrics.}
Table~\ref{tab:rebrand-metrics} summarizes the execution.

\begin{table}[htbp]
\centering
\caption{Client Rebrand Epic Execution}
\label{tab:rebrand-metrics}
\begin{tabular}{ll}
\toprule
\textbf{Metric} & \textbf{Value} \\
\midrule
Total tasks & 12 \\
Max parallelism & 5 concurrent agents (Wave 1) \\
Sequential estimate & 18 hours \\
Parallel estimate & 6.5 hours \\
Speedup factor & 2.8$\times$ \\
Tasks completed & 12 of 12 (100\%) \\
Orchestrator interventions & 3 \\
Test suites / cases & 20 / 174 (0 failures) \\
\bottomrule
\end{tabular}
\end{table}

\needspace{4cm}
\paragraph{Test Fixture Synchronization.}
All three orchestrator interventions involved test assertions falling out
of sync with code changes made by parallel agents. In Wave 1, a pricing
component gained new required props (\texttt{price}, \texttt{period}) that
the existing accessibility test did not supply. In Wave 2, updated navigation
links and SEO strings invalidated tab-order and metadata assertions. None of
these were code defects; they were coordination gaps where one agent's
changes invalidated another agent's test fixtures.

\paragraph{Lesson Learned.}
Structural component changes must propagate to all dependent test fixtures
within the same wave. SPOQ now recommends treating test files as implicit
dependents of the components they exercise, ensuring fixture updates are
included in the same task that modifies the component interface.

\subsection{Failure Modes and Mitigations (Detail)}
\label{subsec:failure-modes-detail}

The following subsections provide detailed analysis of each failure category
identified during SPOQ development and deployment, with specific examples
and comprehensive mitigation strategies.

\subsubsection{Resource Contention Failures}
\label{subsubsec:resource-contention-detail}

\paragraph{Lock File Contention.}
Multiple agents running \texttt{npm install} concurrently cause lock file
conflicts, manifesting as \texttt{EBUSY} errors or corrupted \texttt{node\_modules}
directories. In one deployment, three parallel agents attempted dependency
installation simultaneously, resulting in a 12-minute debugging session.

\textbf{Mitigation:} SPOQ designates a single Wave 0 task for dependency
installation, with subsequent tasks assuming dependencies are available.
Organizations should extend this pattern to any shared mutable resource:
database migrations, cache warming, and artifact generation.

\paragraph{Build Directory Conflicts.}
Concurrent build processes can corrupt shared directories (e.g., \texttt{.next/},
\texttt{dist/}, \texttt{target/}). Symptoms include partial builds, missing
assets, and non-deterministic test failures.

\textbf{Mitigation:} Build verification runs sequentially between waves, not
during waves. For CI/CD integration, consider isolated build directories
per agent or containerized build environments.

\subsubsection{Context and Memory Failures}
\label{subsubsec:context-memory-detail}

\paragraph{Context Window Exhaustion.}
Complex tasks with extensive codebase context can exhaust LLM context windows
(currently 200K tokens for Claude). Symptoms include: agents losing track
of earlier requirements, producing incomplete solutions, or generating
responses that contradict prior instructions. In one knowledge-graph pipeline deployment, an agent
with 180K tokens of context began hallucinating function signatures that
did not exist in the codebase.

\textbf{Mitigation:} Task descriptions should be self-contained, and the
\texttt{files\_to\_touch} specification should limit scope to files the agent
genuinely needs. For large codebases, consider:
\begin{itemize}
    \item Task-specific context manifests listing only relevant files
    \item Summarization of peripheral code rather than full inclusion
    \item Breaking complex tasks into sub-tasks with narrower scope
    \item Context budget monitoring with alerts at 70\% utilization
\end{itemize}

\paragraph{Context Window Management.}
Long QA feedback exhausted agent context windows. The 20-line remediation
limit was introduced to ensure feedback remains actionable without
consuming excessive context. Reviewers now prioritize critical issues
and defer minor suggestions to follow-up tasks.

\subsubsection{Agent Behavioral Failures}
\label{subsubsec:agent-behavioral-detail}

\paragraph{Runaway Detection.}
Agents sometimes enter infinite loops on transient errors, repeatedly
executing the same failing command. In Case Study 1, Task 04 entered a
runaway retry loop, executing \texttt{npm install sonner} over 100 times
before intervention.

\textbf{Mitigation:} SPOQ monitors for repeated identical commands (5+
occurrences) and halts with a human consultation request. Organizations
should implement:
\begin{itemize}
    \item Command deduplication with exponential backoff
    \item Per-task execution time limits (default: 30 minutes)
    \item Anomaly detection for unusual command patterns
\end{itemize}

\paragraph{Validation Gaming.}
Agents optimizing for validation metrics may produce code that passes
automated checks while failing to address underlying requirements. Examples
include: tests that assert \texttt{true === true}, implementations that
short-circuit edge cases, and documentation that restates function signatures
without explaining behavior.

\textbf{Mitigation:} The dual validation framework addresses surface-level
gaming through complementary metrics (test existence vs. test pass rate,
syntactic correctness vs. requirements fidelity). However, sophisticated
gaming requires human oversight:
\begin{itemize}
    \item Periodic human code review of high-complexity tasks
    \item Mutation testing to verify test quality
    \item Requirements traceability audits
    \item Cross-validation between independent agents on critical paths
\end{itemize}

\paragraph{Dependency Resolution Failures.}
Agents may specify incorrect dependency versions, introduce circular
dependencies, or fail to recognize version conflicts. These failures often
manifest only at runtime or during integration testing.

\textbf{Mitigation:}
\begin{itemize}
    \item Lock files (\texttt{package-lock.json}, \texttt{Cargo.lock}) committed
    after Wave 0 dependency installation
    \item Dependency audit as part of code validation metrics
    \item Version pinning policies enforced via task templates
    \item Integration tests run at wave boundaries to detect compatibility issues
\end{itemize}

\needspace{4cm}
\subsubsection{Coordination Failures}
\label{subsubsec:coordination-detail}

\paragraph{Agent Coordination Failures.}
In parallel execution, agents may make conflicting assumptions about shared
state. Example scenarios include:
\begin{itemize}
    \item Two agents independently creating the same utility function with
    different signatures
    \item An agent assuming a database table exists while another agent is
    still creating it
    \item Conflicting CSS class names or component identifiers
\end{itemize}

\textbf{Mitigation:} Wave boundaries serve as synchronization points.
Additional measures include:
\begin{itemize}
    \item Explicit interface contracts defined in task prerequisites
    \item Naming conventions specified in task templates
    \item Linting rules to detect common conflict patterns
    \item Post-wave merge verification before proceeding
\end{itemize}

\paragraph{Human Consultation Latency.}
When agents request human consultation (HaaA), execution pauses until
response. If the human is unavailable, this creates a bottleneck that
can stall entire wave progressions.

\textbf{Mitigation:}
\begin{itemize}
    \item Configurable timeout with fallback behavior (skip, abort, or
    proceed with best-effort)
    \item Escalation to secondary human reviewers
    \item Asynchronous consultation queues for non-blocking requests
    \item Clear documentation of expected response times per request priority
\end{itemize}

\subsubsection{Cost and Resource Failures}
\label{subsubsec:cost-resource-detail}

\paragraph{Cost Runaway Risk.}
Without limits, agents retrying failed tasks or entering verbose output
loops can rapidly consume API credits. A single runaway agent in Case Study 1
consumed approximately \$15 in tokens before detection. Extrapolated to
parallel execution with 10+ agents, uncontrolled failures could exceed
\$100/hour.

\textbf{Mitigation:}
\begin{itemize}
    \item Per-task token budgets (recommended: 50K input, 10K output for
    standard tasks)
    \item Per-epic cost caps with automatic suspension
    \item Real-time cost monitoring dashboards with alerting
    \item Graduated rate limiting: warning at 80\%, throttle at 90\%, halt
    at 100\%
\end{itemize}

\needspace{4cm}
\subsubsection{Security and Isolation Failures}
\label{subsubsec:security-isolation-detail}

\paragraph{Execution Isolation.}
Agents execute commands in the host environment without sandboxing.
Malformed or malicious code could damage the system, exfiltrate data,
or consume resources. While LLM agents are unlikely to be intentionally
malicious, hallucinated commands (e.g., \texttt{rm -rf /}) pose real risks.

\textbf{Mitigation:}
\begin{itemize}
    \item Container-based execution (Docker) with ephemeral workspaces
    \item Read-only filesystem mounts where possible
    \item Network isolation for tasks that do not require external access
    \item Resource limits (CPU, memory, disk quotas)
    \item Command allowlisting for high-risk operations
    \item Audit logging of all executed commands
\end{itemize}

\paragraph{Rollback and Recovery.}
If agent work corrupts the codebase, recovery requires manual git operations.
SPOQ does not currently implement automatic rollback, relying instead on
git history and manual intervention.

\textbf{Mitigation:}
\begin{itemize}
    \item Commit checkpoints at wave boundaries, enabling \texttt{git revert}
    to known-good states
    \item Branch-per-epic isolation with squash merge on completion
    \item Automated backup before epic execution
    \item Recovery runbooks documenting rollback procedures for common failures
\end{itemize}

\subsection{Evaluation Tables}
\label{subsec:evaluation-tables}

\needspace{8cm}
\begin{table}[htbp]
\centering
\caption{Operational Risk Summary}
\label{tab:risk-summary}
\begin{tabular}{lll}
\toprule
\textbf{Risk Category} & \textbf{Detection Signal} & \textbf{Mitigation} \\
\midrule
Context exhaustion & Output truncation & Scope limits, context budgets \\
Cost runaway & Token counter spikes & Per-task budgets with alerts \\
Human bottleneck & Queue timeout & Configurable fallback \\
Lock contention & EBUSY errors & Single-agent dep.\ install \\
Validation gaming & High scores, low quality & Human review, mutation testing \\
Agent conflicts & Merge failures & Interface contracts, wave sync \\
Runaway loops & Repeated commands & Deduplication, time limits \\
Security breach & Unauthorized commands & Container isolation, allowlists \\
Data corruption & Test/build failures & Wave checkpoints, rollback \\
\bottomrule
\end{tabular}
\end{table}

\begin{table}[htbp]
\centering
\caption{Evaluation Summary}
\label{tab:summary}
\begin{tabular}{lccc}
\toprule
\textbf{Metric} & \textbf{UI Epic} & \textbf{Rebrand Epic} & \textbf{Adoption (agg.)} \\
\midrule
Tasks & 13 & 12 & 92 \\
Waves & 2 & 4 & varies \\
Max parallelism & 12 & 5 & 4 \\
Speedup factor & 5.3$\times$ & 2.8$\times$ & 1.3--3.0$\times$ \\
Completion rate & 92\% & 100\% & 100\% \\
Orchestrator interventions & 1 & 3 & varies \\
Rework cycles & 2 & 0 & 0--1 \\
Test cases & --- & 174 & 295 \\
Avg. confidence & --- & 0.92 & 0.90--0.95 \\
\bottomrule
\end{tabular}
\end{table}

\clearpage
\subsection{Cost Analysis and ROI Framework}
\label{subsec:cost-analysis-detail}

A critical consideration for SPOQ adoption is the economic viability of
multi-agent orchestration. We analyze costs under two pricing models
available for Claude API access as of February 2025.

\paragraph{Pricing Model A: Per-Token API.}
Table~\ref{tab:api-pricing} presents the current Claude API pricing structure
and how each tier maps to SPOQ agent roles.

\begin{table}[htbp]
\centering
\caption{Claude API Pricing by Model Tier}
\label{tab:api-pricing}
\begin{tabular}{lccl}
\toprule
\textbf{Model} & \textbf{Input} & \textbf{Output} & \textbf{SPOQ Role} \\
\midrule
Opus 4.6 & \$15/M tokens & \$75/M tokens & Worker agents \\
Sonnet 4.5 & \$3/M tokens & \$15/M tokens & Reviewer agents \\
Haiku 4.5 & \$0.25/M tokens & \$1.25/M tokens & Investigator agents \\
\bottomrule
\end{tabular}
\end{table}

Based on observed token consumption in our case studies, a typical Opus worker
task consumes approximately 25,000 input tokens and 5,000 output tokens,
yielding a per-task cost of approximately \$1.95. For an epic of 13 tasks
(similar to our UI improvements study), the total worker cost reaches
approximately \$28, excluding reviewer and investigator overhead.

\paragraph{Pricing Model B: Flat-Rate Max Plan.}
Anthropic's Max plan (\$200/month) provides 20$\times$ the usage allowance
of a standard Claude Pro subscription. There is no platform-imposed limit on
concurrent Claude Code instances; the practical ceiling is human attention.
Crucially, usage is metered in two separate buckets: \emph{Opus} consumption
and \emph{non-Opus} consumption (Sonnet, Haiku). This two-bucket structure
makes SPOQ's three-tier agent hierarchy an economic optimization as well as
a capability one; reserving Opus tokens for task execution while routing
validation and triage through Sonnet and Haiku preserves Opus headroom for
the work that demands it most.

\begin{itemize}
    \item \textbf{Fixed monthly cost} with predictable budgeting
    \item \textbf{Concurrency:} unlimited instances; bounded only by the
    two usage buckets and human supervisory bandwidth
    \item \textbf{Tiered metering:} Opus budget for Workers; Sonnet/Haiku
    budget for Reviewers and Investigators
    \item \textbf{Daily capacity:} 50--100 tasks (assuming 4--6 task cycles
    across active instances)
    \item \textbf{Effective per-task cost:} \$0.10 at scale (100 tasks/day
    $\times$ 20 working days)
\end{itemize}

At scale, the Max plan reduces per-task costs by approximately 20$\times$
compared to per-token pricing, making aggressive parallelization economically
viable on a single subscription.

\clearpage
\paragraph{The Director Model.}
We propose a scaling paradigm we term the \emph{Director Model}, wherein a
single human engineer orchestrates multiple Claude Code instances under one
Max license:

\begin{itemize}
    \item \textbf{1 planning instance:} Assists with epic decomposition and
    dependency analysis
    \item \textbf{5 execution instances:} Each runs SPOQ's wave-based dispatch
    with parallel sub-agents
    \item \textbf{Daily output:} 50--100 completed tasks
\end{itemize}

In our experience, six concurrent instances represents the human
multitasking limit; the usage buckets are rarely exhausted even at this
level, leaving headroom for sustained parallel execution. This configuration
enables a single engineer
to achieve throughput equivalent to 5--8 traditional engineers, representing
a qualitative shift in individual productivity potential.

\paragraph{ROI Framework.}
Table~\ref{tab:roi} presents a monthly ROI calculation for the Director Model
at typical utilization.

\begin{table}[htbp]
\centering
\caption{Director Model Monthly ROI Estimate}
\label{tab:roi}
\begin{tabular}{lr}
\toprule
\textbf{Component} & \textbf{Value} \\
\midrule
Monthly tasks completed & 2,000 (100/day $\times$ 20 days) \\
Equivalent engineer output & 5--8 engineers \\
Traditional cost (8 engineers @ \$12,500/mo) & \$100,000 \\
Director Model cost (1 Max license) & \$200 \\
Engineer salary (director) & \$12,500 \\
Total Director Model cost & \$12,700 \\
\midrule
\textbf{Monthly savings} & \$87,300 \\
\textbf{ROI multiplier} & 7.9$\times$ \\
\bottomrule
\end{tabular}
\end{table}

Under favorable assumptions (consistent utilization, stable task completion),
the Director Model yields approximately 8$\times$ cost efficiency compared
to traditional staffing.

\paragraph{Caveats and Variability.}
These estimates carry significant uncertainty:

\begin{enumerate}
    \item \textbf{Task complexity variance:} Our cost-per-task figures derive
    from 1--4 hour tasks. Complex tasks requiring extended context windows or
    multiple rework cycles can cost 3--5$\times$ the baseline.

    \item \textbf{Rework overhead:} Failed validations requiring remediation
    add approximately 40\% overhead on affected tasks.

    \item \textbf{Output quality equivalence:} The ROI comparison assumes
    agent-completed tasks are comparable in quality to human-authored work.
    This equivalence has not been empirically validated; agent output may
    require additional review or refinement that narrows the cost gap.

    \item \textbf{Human supervision cost:} The Director Model assumes skilled
    engineers capable of effective agent orchestration. Training and context-switching
    costs are not included.

    \item \textbf{API rate limits:} Per-token pricing may encounter rate limits
    at high parallelism, while Max plan throughput is bounded by
    the usage bucket allowances under sustained heavy parallelism.

    \item \textbf{Quality vs. speed tradeoffs:} Higher parallelism may introduce
    integration issues that offset time savings with debugging overhead.
\end{enumerate}

Despite these caveats, the order-of-magnitude cost advantage suggests that
multi-agent orchestration represents an economically viable paradigm for
software development at scale.

\needspace{6cm}
\subsection{When to Use SPOQ: Decision Framework}
\label{subsec:usage-guidance-detail}

SPOQ's structured approach introduces upfront overhead (task
decomposition, YAML specification, dependency mapping, and validation
gates) that pays off for certain project types but hinders others.
This subsection provides decision guidance for practitioners.

\paragraph{Appropriate Use Cases.}
SPOQ adds value when:

\begin{itemize}
    \item \textbf{Structured features:} New functionality with clear
    requirements spanning multiple components. When you can articulate
    what needs to be built before starting, SPOQ's planning investment
    translates into parallel execution benefits.

    \item \textbf{Multi-file refactors:} Coordinated changes across a
    codebase where parallelism accelerates delivery. Renaming a widely-used
    interface, migrating to a new library, or restructuring module
    boundaries benefit from explicit dependency tracking.

    \item \textbf{Infrastructure work:} Terraform configurations, CI/CD
    pipelines, and deployment automation have well-defined outputs
    and benefit from explicit dependency specification.

    \item \textbf{Team handoffs:} Projects requiring audit trails,
    explainability for stakeholders, or handoffs between engineers
    benefit from SPOQ's journal tracking and task documentation.

    \item \textbf{Scale thresholds:} Our empirical data suggests SPOQ
    provides net benefit for epics with 5+ tasks spanning 4+ hours
    of estimated work. Below this threshold, orchestration overhead
    often exceeds execution time.
\end{itemize}

\paragraph{Inappropriate Use Cases.}
SPOQ's overhead exceeds its value when:

\begin{itemize}
    \item \textbf{Exploratory prototyping:} When requirements are
    unclear and the goal is discovery, SPOQ's upfront planning becomes
    a burden. Rapid iteration with a single agent or direct coding
    serves exploration better.

    \item \textbf{Small fixes:} Single-file bug fixes or minor changes
    where task decomposition time exceeds execution time. A 10-minute
    fix should not require 30 minutes of specification.

    \item \textbf{Creative flow states:} Sessions where the developer
    seeks direct engagement with code, what practitioners sometimes
    call ``vibe coding'', benefit from immediacy, not delegation.
    SPOQ assumes you want to orchestrate, not implement.

    \item \textbf{Urgent hotfixes:} Time-critical patches where
    validation gates delay resolution. When production is down,
    skip the methodology and fix the problem directly.

    \item \textbf{Learning exercises:} When the goal is understanding
    a codebase or technology, doing the work yourself provides
    education that delegation cannot.
\end{itemize}

\paragraph{Decision Framework.}
A simple heuristic: if specifying tasks and dependencies would take
longer than executing the work directly, SPOQ is likely overkill.
Conversely, if you can identify 5+ subtasks with clear deliverables
and the total estimated effort exceeds 4 hours, SPOQ's parallelism
and quality gates typically justify the planning investment.

Table~\ref{tab:when-to-use} summarizes common scenarios:

\begin{table}[htbp]
\centering
\caption{SPOQ Applicability by Project Type}
\label{tab:when-to-use}
\footnotesize
\begin{tabular}{lcp{6.5cm}}
\toprule
\textbf{Project Type} & \textbf{Use SPOQ?} & \textbf{Rationale} \\
\midrule
New feature (5+ files) & Yes & Parallelism benefits; quality gates prevent rework \\
Single-file bug fix & No & Overhead exceeds benefit \\
Exploratory prototype & No & Requirements unclear; planning premature \\
Migration/refactor & Yes & Coordination critical; dependency tracking valuable \\
Learning exercise & No & Flow state preferred; education through doing \\
Production hotfix & No & Urgency trumps process \\
Infrastructure epic & Yes & Well-defined outputs; explicit dependencies \\
Documentation update & Depends & Multi-file: yes; single page: no \\
\bottomrule
\end{tabular}
\end{table}

\paragraph{Overhead vs.\ Payoff.}
SPOQ's overhead includes: (1) epic planning and task decomposition
(30--90 minutes), (2) YAML specification (5--10 minutes per task),
(3) epic validation iteration (10--30 minutes), and (4) journal
tracking and review (ongoing). For a 10-task epic, expect 2--3 hours
of orchestration overhead.

The payoff comes from: (1) parallel execution reducing wall-clock time,
(2) validation gates catching issues before they cascade, (3) reduced
context-switching as agents handle implementation details, and
(4) audit trails simplifying review and handoffs. Across our case studies
and adoption survey, speedups ranged from 1.3$\times$ to 5.3$\times$
depending on dependency structure and parallelization potential.

The break-even point varies by developer productivity and agent
capability, but our experience suggests epics under 5 tasks rarely
justify SPOQ's overhead, while epics over 15 tasks almost always do.
The 5--15 task range requires judgment about parallelization potential
and team coordination needs.
